\section{Knowledge Benchmarks}

知识类 benchmark 是最常见的一类。它们通常采用多选题，测试的是模型是否掌握某些知识，而不一定是真正意义上的``语言理解''。

\subsection{MMLU}

MMLU 包含 57 个学科，覆盖数学、历史、法律、道德等内容。它原本是为 few-shot prompting 时代设计的，用来粗略衡量模型在广泛学科上的知识覆盖。

\begin{itemize}
    \item 优点：覆盖面广，形式标准化。
    \item 缺点：题目质量参差不齐，且越来越容易被过拟合或污染。
    \item 本质：更像知识测试，不是纯语言理解测试。
\end{itemize}

\subsection{MMLU 的具体题目示例与评分机制}

为了让 MMLU 从抽象概念变成具体认知，下面展示一个典型题目（来自 MMLU 的 ``abstract algebra'' 子集）：

\begin{lstlisting}[caption={MMLU 典型题目格式（abstract algebra 子集）},language={}]
Question: Find the degree for the given field
extension Q(sqrt(2), sqrt(3), sqrt(18))
over Q.

A. 0
B. 4
C. 2
D. 6

Answer: B
\end{lstlisting}

评分方式有两种主流做法：

\begin{enumerate}
    \item \textbf{直接概率比较（原始方式）}：不让模型生成文本，而是比较模型在位置上分别给 A/B/C/D 四个 token 的 log-probability，选概率最高的那个作为模型答案。这种方式不需要解析模型输出，但假设模型把答案``集中''在一个 token 上。
    \item \textbf{生成式评估（chain-of-thought 方式）}：让模型先输出推理过程，最后提取答案字母。这种方式更接近真实使用，但需要正则匹配来抽取最终答案，也更容易受 prompt 格式影响。
\end{enumerate}

\begin{warningbox}{同一道题，两种评分方式可能给出不同答案}
直接概率比较和生成式评估的结果不一定一致。有些模型在 log-prob 上选对了，但让它生成推理过程反而会选错（反之亦然）。因此报告 MMLU 分数时，必须注明使用的评分方式。
\end{warningbox}

\begin{knowledgebox}{为什么 MMLU 这类多选题格式有根本性局限}
多选题把所有可能的答案限定在 4 个选项中，这意味着：(1) 随机猜测就有 25\% 的基线；(2) 模型可以通过排除法而非真正理解来答题；(3) 真实场景中用户几乎不会给模型提供选项——他们期望开放式回答。更根本的是，多选格式把``知道答案''和``能生成答案''混为一谈。一个模型可能在看到正确答案时认出它，但如果要从零生成，它可能完全无法给出正确答案。这就像认脸和画脸是完全不同的能力。
\end{knowledgebox}

\subsection{MMLU-Pro}

MMLU-Pro 做了几件事来提升难度：

\begin{itemize}
    \item 去掉了一些噪声和过于简单的题。
    \item 把 4 选项扩成 10 选项。
    \item 用 chain-of-thought 评估，让模型多一点发挥空间。
\end{itemize}

它的意义在于提醒我们：一个基准一旦变得过于饱和，就不再能很好地区分模型。

\subsection{GPQA 和 Humanity's Last Exam}

GPQA 强调的是``Google-proof''，问题由博士级别参与者设计。Humanity's Last Exam 则更进一步，试图构造更难、更多模态、更多领域的题目，并通过多轮审核和 frontier model 过滤来降低平庸题目的比例。

\begin{importantbox}{基准为何会膨胀}
基准题越容易，模型就越快饱和；题越难、越杂、越接近人类研究前沿，越能继续拉开模型差距。但难度提升的同时，构题成本、审核成本和题目质量问题也同步上升。
\end{importantbox}

